\documentclass[11pt]{article}
\usepackage[margin=1in]{geometry}
\usepackage{amsmath,amssymb,amsthm}
\usepackage{booktabs}
\usepackage{hyperref}
\usepackage{xcolor}

\newtheorem{theorem}{Theorem}[section]
\newtheorem{proposition}[theorem]{Proposition}
\newtheorem{lemma}[theorem]{Lemma}
\newtheorem{conjecture}[theorem]{Conjecture}
\theoremstyle{definition}
\newtheorem{definition}[theorem]{Definition}

\newcommand{\RR}{\mathbb R}
\newcommand{\ZZ}{\mathbb Z}
\newcommand{\OO}{O}
\newcommand{\interl}{\mathrel{\preceq}}

\title{Ordered Set Partitions Counted by Big Blocks:\\
Exact Recurrences and Certified Computational Evidence}
\author{Khaled Abdelli}
\date{27 November 2026}

\begin{document}
\maketitle

\begin{abstract}
We study Per Alexandersson's ordered big-block conjecture.  We derive an
integer convolution recurrence, a universal differential recurrence, a
Riccati differential equation, a weighted root-separation reduction, and an
exact obstruction to a naive compatible-window induction.
We describe rigorous exact
methods for deciding real-rootedness and consecutive interlacing, and record
the independently checked parameter ranges.  The distinction between ordering
all blocks and ordering only the large blocks is made explicit.  Finite
certificates are not presented as a proof of the all-parameter conjecture.
\end{abstract}

\section{Definitions and the conjecture}

For \(j\ge2\), let \(\pi=(B_1\mid\cdots\mid B_k)\) be an ordered set partition
of \([n]\), and put
\[
\OO_{n,j}(t)=\sum_{\pi}t^{\#\{i:|B_i|\ge j\}}.
\]
The empty ordered partition contributes \(\OO_{0,j}(t)=1\).  Define
\[
A_j(z)=\sum_{r=1}^{j-1}\frac{z^r}{r!},\qquad
C_j(z)=e^z-1-A_j(z),\qquad
B_j(z,t)=A_j(z)+tC_j(z).
\]
The labelled-sequence construction gives
\[
F_j(z,t):=\sum_{n\ge0}\OO_{n,j}(t)\frac{z^n}{n!}
=\sum_{k\ge0}B_j(z,t)^k
=\frac1{1-B_j(z,t)}.
\]

\begin{conjecture}[Alexandersson, 2026]
For every \(j\ge2\) and \(n\ge0\), the polynomial \(\OO_{n,j}(t)\) is
real-rooted, and \(\OO_{n,j}\interl\OO_{n+1,j}\).
\end{conjecture}

\section{Exact coefficient identities}

\begin{proposition}
The coefficients of \(\OO_{n,j}\) are nonnegative integers and
\(\deg\OO_{n,j}=\lfloor n/j\rfloor\).  Moreover,
\[
\OO_{0,j}=1,\qquad
\OO_{n,j}=
\sum_{r=1}^{j-1}\binom nr\OO_{n-r,j}
+t\sum_{r=j}^{n}\binom nr\OO_{n-r,j}.
\]
\end{proposition}

\begin{proof}
Choose the first block \(B_1\).  If \(|B_1|=r\), its labels can be selected in
\(\binom nr\) ways, and the remaining ordered partition contributes
\(\OO_{n-r,j}\).  The first sum handles unweighted blocks of size less than
\(j\); the second handles a block contributing one factor of \(t\).
Nonnegativity is immediate.  A term \(t^m\) requires at least \(jm\) labels,
and \(m=\lfloor n/j\rfloor\) is attained by using only large blocks.
\end{proof}

For \(j=3\), the recurrence gives
\[
1,\quad 1,\quad 3,\quad 12+t,\quad 66+9t,\quad
450+91t,\quad 3690+973t+20t^2.
\]
At \(t=1\), all block weights are equal, so
\[
F_j(z,1)=\frac1{2-e^z},\qquad
\OO_{n,j}(1)=\sum_{k=0}^n k!{n\brace k},
\]
independently of \(j\).

\section{Riccati and convolution identities}

Write
\[
H_{n,j}(t)=n![z^n]F_j(z,t)^2
=\sum_{r=0}^n\binom nr\OO_{r,j}(t)\OO_{n-r,j}(t).
\]
Since \(F_j'=B_j'F_j^2\), direct coefficient comparison in
\[
B_j'=2+(t-1)\frac{z^{j-1}}{(j-1)!}-(1-B_j)
\]
gives the following exact recurrence.

\begin{theorem}
For \(n\ge0\),
\[
\OO_{n+1,j}=2H_{n,j}-\OO_{n,j}
+(t-1)\binom{n}{j-1}H_{n-j+1,j},
\]
where \(H_{m,j}=0\) for \(m<0\).
\end{theorem}

The final term is signed.  An equivalent form exposes nonnegative structural
coefficients:
\[
F_j'=(B_j+B_j'-1)F_j^2+F_j.
\]
Indeed, if \(b_r=1\) for \(r<j\) and \(b_r=t\) for \(r\ge j\), then the
coefficient of \(z^q/q!\) in \(B_j+B_j'-1\) is \(0\) for \(q=0\),
\(2\) for \(1\le q<j-1\), \(1+t\) for \(q=j-1\), and
\(2t\) for \(q\ge j\).

\section{A universal short recurrence}

Put \(T=t\,\partial_t\). A second exact identity holds for every \(j\ge2\):
\[
\begin{aligned}
\OO_{n,j}={}&n\OO_{n-1,j}+2T\OO_{n-1,j}
+\sum_{r=2}^{j-1}\binom nr\OO_{n-r,j}\\
&+\binom{n-1}{j-1}t\bigl[1+(t-1)\partial_t\bigr]\OO_{n-j,j}.
\end{aligned}
\]
Negative-index terms are omitted. For \(j=3\), this becomes
\[
\OO_{n,3}=n\OO_{n-1,3}+2t\OO'_{n-1,3}
+\binom n2\OO_{n-2,3}
+\frac{t(n-1)(n-2)}2
\bigl[1+(t-1)\partial_t\bigr]\OO_{n-3,3}.
\]

To prove the general identity, write \(w_j=z^{j-1}/(j-1)!\) and \(C_j=e^z-E_{j-1}\).
Since \(C_j'=C_j+w_j\) and \(\partial_tF_j=C_jF_j^2\), differentiation of
\[
tC_jF_j-2\int_0^z TF_j(u,t)\,du
-\frac{t}{(j-1)!}\int_0^z u^{j-1}
\bigl[F_j+(t-1)\partial_tF_j\bigr](u,t)\,du
\]
shows that it is constant in \(z\); its value at \(z=0\) is zero. Extracting
\(n![z^n]\) gives
\[
t\sum_{r=j}^n\binom nr\OO_{n-r,j}
=2T\OO_{n-1,j}
+\binom{n-1}{j-1}t\bigl[1+(t-1)\partial_t\bigr]\OO_{n-j,j}.
\]
Substitution into the first-block recurrence proves the formula.

The same coefficient family, in the snail-race interpretation, appears in
Chen--Zhu, \emph{A Snail Race Problem} (2026), Theorem 3 and Theorem 7.

\section{An inner/outer decomposition}

Let
\[
D_j(z)=\frac1{1-A_j(z)},\qquad g_j(z)=\frac{C_j(z)}{1-A_j(z)},\qquad
\mathcal H_j(z,t)=\frac1{1-tg_j(z)}.
\]
Then
\[
F_j(z,t)=D_j(z)\mathcal H_j(z,t).
\]
Writing \(d_k=[z^k]D_j\) and
\(H_{r,j}(t)=r![z^r]\mathcal H_j(z,t)\), we obtain the exact convolution
\[
\OO_{n,j}(t)=n!\sum_{r=0}^n d_{n-r}\frac{H_{r,j}(t)}{r!}.
\]
Equivalently, with \(N_j=1-A_j\) and \(p_j=z^{j-1}/(j-1)!\),
\[
\OO_{n,j}=H_{n,j}+\sum_{r=1}^{j-1}\binom nr\OO_{n-r,j}. \tag{T1}
\]

The inner family satisfies \(H_{0,j}=1\), \(H_{n,j}=0\) for
\(1\le n<j\), and \(H_{n,j}(0)=0\) for \(n\ge1\). It also satisfies the
linear PDE
\[
N_j\,\partial_z\mathcal H_j
=t\bigl[2+(t-1)p_j\bigr]\partial_t\mathcal H_j+t p_j\mathcal H_j.
\]
Thus the original problem splits into inner real-rootedness of
\((H_{n,j})\) and preservation of real-rootedness/interlacing by the specific
outer transform (T1). The kernels \(g_j\) and \(D_j\) fail the standard
Pólya-frequency hypotheses, so neither reduction is automatic.

\section{A positive coefficient-minor recurrence}

Write \(a_{m,k}=[t^k]\OO_{m,j}\) and
\[
\Delta^{(d)}_{m,k}=a_{m,k+1}a_{m-d,k}-a_{m,k}a_{m-d,k+1}.
\]
The coefficient form of the universal recurrence gives the exact identity
\[
\begin{aligned}
\Delta^{(j)}_{n,k}={}&(n+2k)\Delta^{(j-1)}_{n-1,k}
+2a_{n-1,k+1}a_{n-j,k}\\
&+\sum_{r=2}^{j-1}\binom nr\Delta^{(j-r)}_{n-r,k}\\
&+\binom{n-1}{j-1}\Bigl((k+1)q_k^2+q_kq_{k+1}
-kq_{k-1}q_{k+1}\Bigr),
\end{aligned}
\]
where \(q_k=a_{n-j,k}\), and the first-block recurrence supplies
\(a_{n-1,k+1}\ge\binom{n-1}{j-1}q_{k+1}\). If \(q\) is real-rooted, its
coefficients are log-concave, and the final line is strictly positive.
Thus positive lower-lag minors imply positivity of the lag-\(j\) minor.
This is a genuine positive recurrence, but it controls coefficient ratios,
not the root residues required by (WRS).

\section{A weighted root-separation reduction}

Fix \(n\ge j\), and suppose the window
\(\OO_{n-j,j},\ldots,\OO_{n-1,j}\) is commonly compatible. Put
\[
q=\OO_{n-j,j},\quad p=\OO_{n-1,j},\quad
S_n=\sum_{r=2}^{j-1}\binom nr\OO_{n-r,j},\quad
c=\binom{n-1}{j-1}.
\]
For every root \(\mu\) of \(q\), impose the weighted root-separation inequality
\[
-\frac{n p(\mu)+2\mu p'(\mu)+S_n(\mu)}{q'(\mu)}
>c\mu(\mu-1). \tag{WRS}
\]
Then \(\OO_{n,j}\) has \(\deg q+1\) distinct negative real roots and
interlaces \(q\).

Indeed, at a root \(\mu\) of \(q\),
\[
t[1+(t-1)\partial_t]q\big|_{t=\mu}=\mu(\mu-1)q'(\mu).
\]
Thus (WRS) is exactly the assertion
\(\operatorname{sgn}\OO_{n,j}(\mu)=-\operatorname{sgn}q'(\mu)\).
These signs alternate at the roots of \(q\). Since
\(\deg\OO_{n,j}=\deg q+1\), the signs at \(-\infty\), between consecutive
roots of \(q\), and at \(0\) force \(\deg q+1\) distinct roots by the
intermediate value theorem.

Mere compatibility is insufficient. For \(j=2,n=5\), take
\(p=(t+2)(t+4)\) and \(q=t+3\). Then \(q\) interlaces \(p\), and also
interlaces \(R_5(p)=9t^2+42t+40\). Nevertheless, at \(\mu=-3\), the left side
of (WRS) is \(5\) and the right side is \(48\). The resulting test polynomial
\(17t^2+50t+40\) has discriminant \(-220\). A successful induction must
therefore prove the quantitative inequality (WRS), not only a compatible
window.

The Bézout-matrix formulation gives the same gap. For the root-positive
Bézoutian \(M(q,p)\), bilinearity gives the all-index identity
\[
M(q,\OO_n)=M(q,A_n)+c\,v(q)v(q)^T+c\,M(q,t(t-1)q'),
\]
where \(v(q)\) is the coefficient vector of \(q\). At every root \(\mu\) of
\(q\), the intrinsic kernel has diagonal value
\(-\mu(\mu-1)q'(\mu)^2<0\). Thus positive definiteness of \(M(q,A_n)\)
alone cannot suffice; the actual rows must quantitatively dominate this
negative kernel.

Equivalently, at a root \(\mu\) of \(q\), the first-block recurrence writes
\[
\frac{\OO_n(\mu)}{q'(\mu)}
=\sum_{r=1}^{j-1}\binom nr\frac{\OO_{n-r}(\mu)}{q'(\mu)}
+\mu\sum_{r=j+1}^{n}\binom nr\frac{\OO_{n-r}(\mu)}{q'(\mu)}.
\]
Recent-window compatibility controls the first sum, but older rows can change
sign at old roots. The missing all-index input is a quantitative domination
bound for this entire older tail.

\section{A matrix-route obstruction}

Write \(a_{n,m}=[t^m]\OO_{n,j}\), \(D_j=(1-A_j)^{-1}\), and
\(h_j=(C_j/z^j)D_j\). Then
\[
\frac{a_{n,m}}{n!}=R^{(j)}_{n-jm,m},\qquad
R^{(j)}_{q,m}=[z^q]D_j(z)h_j(z)^m.
\]
Thus \(R^{(j)}\) is a proper Riordan array, but the target is a column shear.
General shear preservation is false: there is a proper Riordan array \(R\)
with a totally nonnegative production matrix for which
\[
\det R[\{4,6\},\{0,1\}]
=\det\begin{pmatrix}14&28\\3&132\end{pmatrix}=-4.
\]
The target therefore requires a special sheared-minor theorem or a planar
network for its weighted block-size paths; ordinary total positivity of
\(R^{(j)}\) is not enough. The natural weighted path digraph is also
nonplanar; planarizing it yields a reflected, sign-twisted determinant
inequality rather than the desired minors.

\section{Related ordered lifts are different}

Let \(P_{n,j}\) count ordinary set partitions by their large blocks.  Its EGF
is \(e^{B_j}\).  Ordering only the \(m\) large blocks replaces \(t^m\) by
\(m!t^m\), and therefore has EGF
\[
\frac{e^{A_j(z)}}{1-tC_j(z)}.
\]
This is \emph{not} \(1/(1-A_j-tC_j)\), which orders all blocks.  The
distinction is essential when applying theorems for restricted big-block
families to the present conjecture.

\section{Rigorous finite certification}

All coefficient generation uses integer arithmetic.  Root isolation uses
FLINT's integer-polynomial routine with Arb complex balls.  A root is
certified non-real when its imaginary ball excludes zero; a real-rooted
verdict requires every returned root enclosure to lie on the real axis.
Consecutive interlacing is checked by comparing rigorous lower and upper
endpoints of ordered real-root balls, including every degree-increasing
transition.

Independent checks use exact Sturm sequences, exact discriminants, and
Newton's ultra-log-concavity inequalities.  A negative discriminant is an
immediate obstruction; a zero discriminant is not, since repeated real roots
are allowed.  Floating-point root finders are never used as certificates.

% Draft placeholder. Replaced only after the independent audits report.
\section{Computational results}
The computational result tables are being generated from the validated
JSONL audit logs.


\section{Conclusion and remaining work}
The all-parameter conjecture is not inferred from finite computation. The
remaining consecutive-interlacing target is the root-sign inequality
\[
\frac{\OO_{n+1,j}(\mu)}{\OO_{n,j}'(\mu)}<0
\]
at every root \(\mu\) of \(\OO_{n,j}\). The lag-\(j\) WRS route is equivalently
the negative-axis Wronskian inequality
\(W(\OO_{n,j},\OO_{n-j,j})(t)>0\), with the full older tail isolated as
\(W(H_{n,j},\OO_{n-j,j})\). Standard PF, CZDS, Favard, Riordan-shear,
planar-network, crossing-involution, threshold-renewal, and Bézout-positivity
routes all have exact obstructions. Fixed-\(K\) edge asymptotics also do not
control the \(\Theta(m)\) middle zeros. The certified computations identify
the precise unchecked frontier.

\begin{thebibliography}{9}
\bibitem{Ale26}
P. Alexandersson.
\newblock Open problems on real-rootedness.
\newblock SymCat, 2026.
\newblock \url{https://www.symmetricfunctions.com/realRootedOpenProblems.htm#conjOrderedBigBlocks}.

\bibitem{ChenZhu26}
Y. Chen and W. Zhu.
\newblock A Snail Race Problem.
\newblock arXiv:2607.02085, 2026.
\newblock \url{https://arxiv.org/abs/2607.02085}.

\bibitem{ChenLW16}
X. Chen, H. Liang, and Y. Wang.
\newblock Total positivity of Riordan arrays.
\newblock arXiv:1601.05637, 2016.

\bibitem{ChenFR25}
X. Chen, L. Fu, and J. Ruan.
\newblock A unified planar network approach to total positivity of combinatorial matrices and real-rootedness of polynomials.
\newblock arXiv:2512.08369, 2025.

\bibitem{Issue108}
\newblock Public computational audit, Issue \#108, 2026.
\newblock \url{https://github.com/1337isnot1337/solved-conjectures/issues/108}.
\end{thebibliography}

\end{document}
